\section{Substrate model: lattice microphysics and continuum (IR) effective description}
\label{sec:continuum-model}

\subsection{Ontology and kinematics}
\label{subsec:ontology-kinematics}

The fundamental object of this theory is a \emph{4D brane lattice embedded in a 4D Euclidean ambient with codimension zero}: the lattice is the substrate, and there is no extra perpendicular direction in the ambient.
The ambient $\R^4$ is fully symmetric --- no preferred direction --- and serves only as a stage on which Euclidean distances are measured.
All asymmetries that we will encounter (the distinction between time and the three spatial directions, the split between gauge and gravity sectors) are properties of the brane and of the inside observer's frame, not of the ambient.

Let $\Omega\subset\R^4$ be the intrinsic (material) brane domain (either a periodic cell $\Omega=\T^4$ or a bounded hypercube with clamped boundary).
The fundamental field is the embedding map
\begin{equation}
  \vecX(x)\in\R^4,
  \qquad x=(x^1,x^2,x^3,x^4)\in\Omega,
  \label{eq:fundamental-field}
\end{equation}
whose components $X^A(x)$ ($A=1,\dots,4$) give the brane configuration in the ambient.
The four ambient directions are geometrically equivalent.
The asymmetry that picks one of them as the \emph{timelike} direction is a property of the brane action $S[\vecX]$, not of the ambient: the action has Lorentzian sign structure --- the standard $T-V$ Lagrangian, recast on the lattice as link-energy terms with opposite sign for the timelike-direction contributions.
Globally, cosmological boundary conditions (a Big Bang vertex emanating along one specific ambient direction) fix this choice unambiguously, so every inside observer agrees on which axis is time.

\paragraph{Working spacelike-slice description.}
For technical work that does not require explicit 4D treatment --- the linear-regime dispersion of \Cref{sec:emergent-relativity}, the gauge-sector diagnostics of \Cref{sec:geophase}, and the localized-mode analysis of \Cref{sec:localized} --- it is convenient to slice the brane perpendicular to the timelike direction.
Each spacelike slice is parameterised by $x=(x^1,x^2,x^3)\in\Omega_3\subset\R^3$, with the timelike coordinate $t:=x^4$ acting as the slicing parameter.
In this description the embedding takes the equivalent form
\begin{equation}
  \vecX(x,t)\in\R^4,
  \qquad x=(x^1,x^2,x^3)\in\Omega_3,\quad t\in\R,
  \label{eq:fundamental-field-sliced}
\end{equation}
that of a three-dimensional material domain embedded in the 4D ambient.
The fourth ambient component, denoted $u$ in this slicing, is the inside observer's chosen timelike component of $\vecX$; relative to the spacelike slice it plays the role of a transverse ``amplitude'' direction.
With this identification, all the constitutive, kinematic, and gauge-sector developments below carry through verbatim, now understood as describing the inside observer's spacelike slice.

\subsection{Discrete lattice substrate (microphysical postulate)}
\label{subsec:lattice-substrate}

The continuum model is treated as the long-wavelength effective theory of a \emph{4D cubic lattice} substrate.
Lattice sites are labeled by $n\in\mathbb{Z}^4$ with spacing $a$ and reference positions $x_n=a\,n$; each site carries an embedded configuration $\vecX_n\in\R^4$.
The discrete brane action is a sum of link-energy contributions with the Lorentzian sign structure of \Cref{subsec:ontology-kinematics}: contributions involving the timelike direction enter with opposite sign from those involving the three spacelike directions.

On a fixed spacelike slice (perpendicular to the timelike direction), sites are labeled by $n\in\mathbb{Z}^3$ with embedded configurations $\vecX_n(t)\in\R^4$ where $t$ is the slicing parameter.
For a spacelike neighbour stencil $\mathcal N\subset\mathbb{Z}^3$, the slice-restricted potential energy reads
\begin{equation}
  U_{\mathrm{lat}}[\{\vecX_n\}]
  =
  \frac12\sum_n\sum_{\delta\in\mathcal N}
  k_{\delta}
  \Big(\,\lvert \vecX_{n+\delta}-\vecX_n\rvert-\alpha a\lvert\delta\rvert\,\Big)^2,
  \label{eq:lattice-energy}
\end{equation}
with stiffnesses $k_{\delta}$ grouped, for example, into axis-aligned, face-diagonal, and body-diagonal shells.
The uniform prestretch parameter $\alpha$ fixes the ratio between the held lattice spacing and the stress-free link
length.
Link contributions along the timelike direction enter the full brane action with opposite sign and are absorbed into the kinetic term $\tfrac12\rho_m|\partial_t\vecX|^2$ of the slice Lagrangian (\Cref{eq:lagrangian-density} below); their explicit lattice form is developed in \Cref{subsec:discrete-4d-action}.

\paragraph{Geometric nonlinearity already present at the microscopic level.}
\Cref{eq:lattice-energy} makes an important modeling point explicit.
Even if the link energy is Hooke-linear in the scalar extension,
its dependence on the node coordinates is nonlinear because the extension is the full Euclidean distance in $\R^4$.
This Pythagorean nonlinearity is the microscopic origin of the amplitude--lateral coupling that later appears in the
continuum as the induced-metric term $\partial_i u\,\partial_j u$.
No additional nonlinear interaction field is required for this effect.

\paragraph{Minimal stencil and lab-frame anisotropy.}
The canonical realisation of the lattice in this paper is the \emph{6-neighbor axial-only} stencil:
each site connects to its six nearest axial neighbours $\pm\hat{e}_i$.
Diagonal-shell (face- and body-diagonal) bonds are intentionally absent.
On this stencil the long-wavelength dispersion is cubically anisotropic at leading order in $|k|a$:
for example, the longitudinal speed along $[100]$ exceeds the longitudinal speed along $[111]$
by a factor $1/\sqrt{(3-2\alpha)/3}$, evaluating to roughly $7.5\%$ at the default operating point $\alpha=0.2$.
This anisotropy is not retuned away; it is the structural feature that the gauge-sector emergence story in
\Cref{sec:geophase,sec:effective-field} ultimately rests on.

\paragraph{Dual-observer framework: lab anisotropy vs internal isotropy.}
The lab (outside) observer knows the lattice and resolves direction-dependent wave speeds.
The internal (soliton) observer is built from substrate excitations and, by hypothesis, has rods, clocks, and signal
cones that all renormalize with the same direction-dependent local wave speed.
Their measured speed of light, defined operationally as the ratio of an emergent rod length to an emergent clock period,
is therefore conjectured to be direction-independent in their own units.
This is a strong assumption: it promotes ``operational isotropy'' from a passive consistency condition to the
load-bearing mechanism by which Lorentz-like kinematics emerge for internal observers despite manifest lab-frame
anisotropy.
The decisive empirical test is whether two solitons in different orientations report the same effective metric;
this is the central numerical objective of \Cref{sec:emergent-relativity} and is operationalised in the validation
roadmap (Sprint~3).

\paragraph{Brillouin zone and geometric-phase diagnostics.}
The lattice spacing $a$ defines a Brillouin zone $k_i\in[-\pi/a,\pi/a]$.
Because Berry curvature is naturally defined over parameter spaces such as $\mathbf{k}$-space eigenbundles,
the lattice formulation provides a direct route to Berry-curvature diagnostics on a discretized momentum lattice
(see \Cref{sec:geophase}).

\subsection{Explicit discrete 4D brane action}
\label{subsec:discrete-4d-action}

The spacelike-slice description of \Cref{subsec:lattice-substrate} treats the timelike direction implicitly, absorbing it into the kinetic term $\tfrac12 m |\dot{R}|^2$.
For the foundational 4D-in-4D ontology of \Cref{subsec:ontology-kinematics} it is essential to make the timelike direction an explicit lattice link, placing all four spatial directions on equal footing as far as the link topology is concerned and making the Lorentzian sign structure a single explicit choice.
Label each node by a spacelike triple $p=(i,j,k)$ and a timelike index $l\in\mathbb{Z}$, with 4D embedded position $R_p^l\in\R^4$.
The six spacelike axial neighbours form the stencil $\mathcal{N}_s=\{\pm\hat{e}_x,\pm\hat{e}_y,\pm\hat{e}_z\}$ and the two temporal neighbours form $\mathcal{N}_t=\{l\pm 1\}$.

\paragraph{The action.}
The slice-restricted potential is the central-force sum over $\mathcal{N}_s$:
\begin{equation}
  V^l = \frac{k_s}{4}\sum_{p}\sum_{\delta\in\mathcal{N}_s}
        \Bigl(\bigl|R_{p+\delta}^l - R_p^l\bigr| - \alpha a\Bigr)^2,
  \label{eq:slice-potential}
\end{equation}
and the temporal kinetic term on the link $l\to l+1$ is
\begin{equation}
  T^{l+\frac{1}{2}} = \sum_{p}\frac{m}{2}
        \left(\frac{R_p^{l+1}-R_p^l}{\Delta t}\right)^{\!2}.
  \label{eq:temporal-kinetic}
\end{equation}
With the Lorentzian sign assignment --- spacelike links enter with $-$ (potential), temporal links enter with $+$ (kinetic); this is Assumption~A5 of the derivation and is \emph{posited}, not derived: it is the property of the brane action that selects one of the four equivalent ambient directions as timelike (backbone~\#21) --- the explicit discrete action is
\begin{equation}
  \boxed{
  S[R] = \sum_l \Delta t\,\Bigl(T^{l+\frac{1}{2}} - V^l\Bigr).
  }
  \label{eq:discrete-4d-action}
\end{equation}
This is a sum over $6+2$ links per node with a single sign flip for the temporal family.

\paragraph{Euler--Lagrange identity: Verlet is the variational integrator.}
Varying \Cref{eq:discrete-4d-action} with respect to an interior node $R_p^l$ (the node appears in $T^{l\pm\frac{1}{2}}$ and $V^l$) yields
\begin{equation}
  m\,\frac{R_p^{l+1} - 2R_p^l + R_p^{l-1}}{\Delta t^2}
  \;=\;
  F_p^l
  \;:=\;
  -\frac{\partial V^l}{\partial R_p^l}
  = k_s\!\sum_{\delta\in\mathcal{N}_s}
    \Bigl(\bigl|R_{p+\delta}^l - R_p^l\bigr| - \alpha a\Bigr)
    \,\widehat{(R_{p+\delta}^l - R_p^l)},
  \label{eq:el-verlet}
\end{equation}
which is precisely the Störmer--Verlet update $R_p^{l+1} = 2R_p^l - R_p^{l-1} + (\Delta t^2/m)\,F_p^l$.
This is the discrete-mechanics result (Marsden--West): the Euler--Lagrange equations of the discrete action $S$ are exactly Störmer--Verlet, and the flow is symplectic.
Therefore \emph{Verlet is the discrete variational integrator of the action \eqref{eq:discrete-4d-action}}.
The forward initial-value problem (prescribe $R^0, R^1$, march $l\to l+1$) and the foundational 4D block boundary-value problem (prescribe spacelike-slice data at $l=0$ and $l=N$, root-find the interior) share \emph{identical local stencil} \eqref{eq:el-verlet} and differ only in their global boundary conditions.

\paragraph{Long-wavelength cone speeds.}
Taking the continuum limit with $\rho_m = m/a^3$ and wave vector $\mathbf{k}=k\hat{e}_x$, the Christoffel matrix is diagonal on the axial stencil $\mathcal{N}_s$ (no off-diagonal coupling at any $\alpha$), giving the lattice light-cone speeds
\begin{equation}
  c_L^2 = \frac{k_s\,a^2}{m},
  \qquad
  c_T^2 = (1-\alpha)\,\frac{k_s\,a^2}{m}.
  \label{eq:cone-speeds}
\end{equation}
In dimensionless units $k_s=a=m=1$ these reduce to $c_L=1$, $c_T=\sqrt{1-\alpha}$.
At the default operating point $\alpha=0.2$ one obtains $c_T/c_L=\sqrt{0.8}\approx 0.894$, a residual $\sim 10.6\%$ longitudinal--transverse gap that is the leading-order consequence of prestretch; it is not retuned away (backbone~\#8, \#15).

\paragraph{Outlook: foundational block solver.}
The Lorentzian action \eqref{eq:discrete-4d-action} is a saddle (unbounded below), so a foundational block solver must root-find $\nabla S=0$ rather than minimize $S$.
The well-posedness of the two-time boundary-value problem and the design choice between the kinetic-increment temporal model (model~a, baseline, used throughout this paper) and a central-force temporal spring model (model~b) are open questions tracked outside the manuscript.

\subsection{Induced metric and isotropic reference metric}

The embedding induces a (time-dependent) Riemannian metric on the brane:
\begin{equation}
  g_{ij}(x,t) := \partial_i \vecX(x,t)\cdot \partial_j \vecX(x,t),
  \qquad i,j=1,2,3.
  \label{eq:induced-metric}
\end{equation}
This definition is invariant under rigid motions $\vecX\mapsto Q\vecX+a$ with $Q\in O(4)$ and $a\in\R^4$.
Rotational invariance within the brane is obtained by constructing the stored energy density from invariants of $g$
relative to an isotropic reference metric.

We choose a homogeneous, isotropic reference metric
\begin{equation}
  g^0_{ij}=\ell_0^2\,\delta_{ij},
  \label{eq:reference-metric}
\end{equation}
where $\ell_0>0$ is the continuum analogue of a discrete spring rest length.
Define the mixed relative metric
\begin{equation}
  \widehat g := (g^0)^{-1}g.
  \label{eq:relative-metric}
\end{equation}
Any isotropic hyperelastic energy density can be written as $W=W(\widehat g)$, equivalently as a function of the
principal stretches or invariants of $\widehat g$.

\paragraph{Continuum analogue of pre-tension.}
The reference metric $g^0$ encodes the stress-free geometric scale.
If boundary conditions enforce a held ground-state scale $h_\star$ (uniform), then
\begin{equation}
  \alpha:=\frac{\ell_0}{h_\star}\in(0,1]
  \label{eq:alpha-def}
\end{equation}
measures the mismatch between stress-free scale and held scale.
For $\alpha<1$ the brane is uniformly pre-stressed (continuum ``pre-tension''). This mismatch controls how strongly
transverse gradients backreact on in-brane stresses through the geometric nonlinearity of $g_{ij}(\vecX)$.
In the lattice substrate picture (\Cref{subsec:lattice-substrate}), the same $\alpha$ multiplies the geometric
neighbor distances to set spring rest lengths; \Cref{eq:alpha-def} is the continuum encoding of that uniform prestretch.

\subsection{Constitutive law: St.\ Venant--Kirchhoff (Green--Lagrange strain)}

As a concrete baseline we use the St.\ Venant--Kirchhoff (StVK) energy expressed in terms of the
Green--Lagrange strain tensor
\begin{equation}
  E := \tfrac12\big(\widehat g - I\big)
  \qquad\text{(with $I$ the $3\times 3$ identity)}.
  \label{eq:green-lagrange-strain}
\end{equation}
The StVK stored energy density is
\begin{equation}
  W_{\text{StVK}}(E)
  = \mu\,\mathrm{tr}(E^2)+\frac{\lambda}{2}\,\big(\mathrm{tr}E\big)^2,
  \label{eq:stvk}
\end{equation}
with Lam\'e parameters $\mu>0$ and $\lambda$ (units: Pa). In terms of Poisson ratio $\nu\in(0,1/2)$,
\begin{equation}
\nu=\frac{\lambda}{2(\lambda+\mu)},
\qquad
\lambda=\frac{2\mu\nu}{1-2\nu}.
\label{eq:poisson-lame}
\end{equation}
This choice is not claimed unique; it is used because it is isotropic, analytic, and captures the
geometric nonlinearity through $g_{ij}(\vecX)$.
It also reduces to standard linear elasticity in the small-strain limit, making the longitudinal/transverse
wave speeds and the linearized branch structure transparent.
We emphasize that StVK is employed here as a \emph{baseline constitutive law} for analytic tractability.
Its relevant limitation is specific rather than a generic large-strain nicety: StVK is neither polyconvex nor coercive in compression --- its stored energy does not diverge as $\det F\to 0^+$, so it neither penalizes local volume collapse and element inversion nor stiffens sufficiently under strong compression.
This feature is structural, not an artifact of the continuum approximation: the microscopic central-force lattice shares it, since the pair energy $\tfrac12 k_\delta(|R_{p+\delta}-R_p|-\alpha a)^2$ stays bounded --- tending to $\tfrac12 k_\delta(\alpha a)^2$ --- as a link collapses, $|R_{p+\delta}-R_p|\to 0$.
The lattice spacing $a$ merely masks the instability (configurations cannot compress below the grid) that the continuum exposes.
For the linear and moderate-strain results of \Cref{sec:emergent-relativity,sec:geophase} this is immaterial; for the strongly localized modes of \Cref{sec:localized} it must be controlled --- either by verifying that the relevant configurations remain non-inverting, or by substituting a polyconvex hyperelastic energy (e.g.\ neo-Hookean, Mooney--Rivlin, Ogden-type) within the same induced-metric framework, which leaves the geometric mechanism of transverse--lateral backreaction unchanged.
The geometric quartic that stabilizes solitons against Derrick (uniform-rescaling) collapse is a distinct safeguard and does not by itself remove this local compression instability.

A second, distinct caveat concerns the \emph{quartic} order. StVK is quadratic in the Green--Lagrange strain $E$, whereas the microscopic central-force spring is quadratic in the \emph{stretch} $(|R_{p+\delta}-R_p|-\alpha a)$; the two coincide at quadratic order (identical wave speeds, \Cref{sec:emergent-relativity}) but differ at quartic order. Writing the exact link energy as $\tfrac12 k_s|R_{p+\delta}-R_p|^2 - k_s\alpha a\,|R_{p+\delta}-R_p| + \text{const}$, the squared-norm term is exactly quadratic and the \emph{entire} anharmonic sector is the Euclidean-norm term $-k_s\alpha a\,|R_{p+\delta}-R_p|$, which is proportional to $\alpha$. The lattice-exact geometric quartic therefore has coefficient $\propto k_s\alpha/a$ (vanishing as $\alpha\to0$), whereas a naive StVK identification locks it to $\mu\propto(1-\alpha)$ and inverts the $\alpha$-scaling. Quantitative soliton sizing (\Cref{sec:localized}) must accordingly use the lattice-derived quartic rather than the StVK one.

\paragraph{Spacelike-slice versus 4D placement.}
The strain \eqref{eq:green-lagrange-strain} is a spacelike-slice object: $E$ is $3\times 3$ and is referred to the Euclidean reference metric of the slice.
Its placement in the foundational 4D action (\Cref{subsec:discrete-4d-action}) admits the two readings of the timelike-link choice.
In the baseline, $W_{\text{StVK}}$ is the potential $V$ and the timelike direction contributes a separate kinetic term, so $S=\int(T-V)$ is ordinary elastodynamics and \eqref{eq:stvk} stands verbatim.
In the symmetric alternative, the strain is promoted to a $4\times 4$ Green--Lagrange tensor $E_{\mu\nu}=\tfrac12(g_{\mu\nu}-\eta_{\mu\nu})$ referred to a Lorentzian reference $\eta_{\mu\nu}$.
In either reading the Lorentzian signature resides in the \emph{sign structure of the action} (\Cref{subsec:ontology-kinematics}), not in the ambient or induced metric: the ambient is Euclidean, so $g_{\mu\nu}=\partial_\mu\vecX\cdot\partial_\nu\vecX$ is positive-definite and the timelike direction is distinguished only by the opposite sign with which its strain enters $S$.

\subsection{Lagrangian, action, and equations of motion}

Let $\rho_m>0$ be the brane mass density (kg/m$^3$). The Lagrangian density is
\begin{equation}
  \Lag(\vecX,\partial_t\vecX,\nabla\vecX)
  =\frac{\rho_m}{2}\,|\partial_t\vecX|^2
  -W\!\left(g(\vecX);g^0\right).
  \label{eq:lagrangian-density}
\end{equation}
The action is
\begin{equation}
  S[\vecX] = \int_{\R}\dd t \int_{\Omega} \dd^3 x\;\Lag.
  \label{eq:action}
\end{equation}

Define the first Piola stress (as a momentum conjugate to $\partial_i X^A$)
\begin{equation}
  P_A^{\;i} := \frac{\partial W}{\partial(\partial_i X^A)}.
  \label{eq:piola-def}
\end{equation}
For energies depending on $g_{ij}$ one can write $P_A^{\;i}=2S^{ij}\partial_j X^A$ with a symmetric
second Piola-type tensor $S^{ij}=\partial W/\partial g_{ij}$.
The Euler--Lagrange equations take divergence form
\begin{equation}
  \rho_m\,\partial_{tt} X^A = \partial_i P_A^{\;i},
  \qquad A=1,\dots,4,
  \label{eq:eom}
\end{equation}
showing explicitly that isotropy (absence of built-in preferred directions) follows from writing
$W$ in terms of invariants of the induced metric relative to the reference metric $g^0$.

\subsection{Spherical material coordinates (coordinate change only)}
\label{subsec:spherical-material-coords}

\paragraph{Important: this is only a coordinate change.}
All dynamical statements remain those of the coordinate-free formulation
$\vecX:\Omega\times\R\to\R^4$ with induced metric $g_{ij}=\partial_i\vecX\cdot\partial_j\vecX$.
We introduce spherical coordinates on the \emph{intrinsic/material} domain as a convenient chart for localized
excitations; this does \emph{not} constrain the embedding $\vecX$.

Let $(r,\theta,\varphi)$ be the standard spherical chart on (a subset of) $\R^3$, with Jacobian
$J(r,\theta,\varphi)=r^2\sin\theta$.
Writing the intrinsic coordinate as $q^a\in\{r,\theta,\varphi\}$ and denoting derivatives in this chart by $\partial_a$,
the equations of motion become
\begin{equation}
  \rho_m\,\partial_{tt}X^A
  =\frac{1}{J(q)}\,\partial_a\!\Big(J(q)\,P_A^{\;a}\Big),
  \qquad A=1,\dots,4,
  \label{eq:eom-spherical}
\end{equation}
which is exactly the Cartesian divergence form \eqref{eq:eom} written with the standard Jacobian factor.

\subsection{Linearization and polarized branches}

Linearize around a homogeneous background embedding $\vecX_\star(x)=(x,0)$ and small displacements
$\vecu(x,t)=\vecX(x,t)-\vecX_\star(x)$.
Under standard regularity assumptions on $W$, one obtains a linear wave system with (in general) multiple
polarization branches, including transverse and longitudinal modes.
For later use, we emphasize two generic features:
(i) different branches may have different characteristic speeds (and may become dispersive at shorter wavelengths),
and (ii) for $\alpha<1$ the pre-stress couples transverse excursions (in the embedding) back into in-brane stresses
through the induced metric \eqref{eq:induced-metric}.
These branches provide the substrate carriers for the geometric-phase and soliton constructions developed below.

\subsection{Transverse--lateral coupling and an emergent ``gravity'' channel}
\label{subsec:gravity-channel}

A distinctive feature of an embedded elastic medium is that transverse gradients change intrinsic distances.
To make this explicit without removing in-brane dynamics, we use a near-identity decomposition around the
homogeneous background embedding $\vecX_\star(x)=(h_\star x,0)$:
\begin{equation}
  \vecX(x,t)=\big(h_\star x^1+\xi^1(x,t),\,h_\star x^2+\xi^2(x,t),\,h_\star x^3+\xi^3(x,t),\,u(x,t)\big),
  \qquad \vec\xi(x,t)\in\R^3.
  \label{eq:near-identity-embedding}
\end{equation}
The induced metric is then
\begin{align}
  g_{ij}
  &= \partial_i\vecX\cdot\partial_j\vecX \nonumber\\
  &= (h_\star\delta_i^{\,a}+\partial_i\xi^a)(h_\star\delta_j^{\,a}+\partial_j\xi^a)+\partial_i u\,\partial_j u \nonumber\\
  &= h_\star^2\delta_{ij}+h_\star(\partial_i\xi_j+\partial_j\xi_i)+\partial_i\vec\xi\cdot\partial_j\vec\xi+\partial_i u\,\partial_j u.
  \label{eq:metric-near-identity}
\end{align}
Setting $\vec\xi\equiv 0$ recovers the familiar graph picture used only for intuition; the ``gravity'' channel,
however, relies on the allowed in-brane relaxation $\vec\xi\neq 0$ sourced by the quadratic term
$\partial_i u\,\partial_j u$.

In the small-slope regime, the in-brane strain has the schematic form
\begin{equation}
  \varepsilon_{ij}\approx \tfrac12(\partial_i\xi_j+\partial_j\xi_i)+\tfrac12\,\partial_i u\,\partial_j u,
  \label{eq:strain-near-identity}
\end{equation}
so the quasi-static balance $\partial_i\sigma_{ij}\approx 0$ with
$\sigma_{ij}=\lambda\,\mathrm{tr}(\varepsilon)\delta_{ij}+2\mu\,\varepsilon_{ij}$ makes it explicit that
the term $\partial_i u\,\partial_j u$ acts as a forcing/eigenstrain-like source for the slow in-brane field
$\vec\xi$.

\paragraph{Pythagorean lengthening and lateral contraction (informal picture).}
For a small material displacement $\dd x$, the ambient line element of the graph embedding satisfies
\begin{equation}
  \dd s^2 = h_\star^2\,|\dd x|^2 + (\nabla u\cdot \dd x)^2.
  \label{eq:line-element-graph}
\end{equation}
Thus any region with nonzero slope $\nabla u$ necessarily \emph{lengthens} intrinsic distances compared to the flat
configuration.
In a discrete spring-lattice realization this is the familiar Pythagorean effect
$|\Delta \vecX|=\sqrt{|\Delta x|^2+(\Delta u)^2}$: transverse variation increases spring length even when lateral
coordinates are unchanged.
In the actual dynamics the in-brane displacement $\vec\xi$ is not constrained; under pre-tension ($\alpha<1$) the
additional metric contribution $\partial_i u\,\partial_j u$ acts as an effective source in the in-brane equilibrium
equations, producing a compensating contraction/dilation pattern rather than a purely local increase of tensile stress.
If the medium is pre-tensioned ($\alpha<1$), the additional length is paid for by increased tensile stress,
which produces net lateral forces that pull material points \emph{toward} the region of transverse excitation.
For a localized bump or rotating transverse pattern, the surrounding brane therefore develops an inward bias
(contraction) rather than an outward pressure.

In this qualitative sense, the ``gravitational mass'' of an excitation is tied to the elastic energy stored in
transverse gradients (and, more generally, to embedded deformation energy), while the ``gravitational field'' is the
slowly varying in-brane strain/dilation pattern sourced by that energy.
The sign is fixed by geometry: increasing transverse arc length increases elastic energy and therefore increases
tension, which pulls laterally inward.

With homogeneous pre-stress ($\alpha<1$), this geometrically induced strain backreacts on the in-brane stress field
through the constitutive law, effectively coupling transverse excitation energy into slow modulations of the
in-brane configuration.

\paragraph{Unification across all four directions.}
The Pythagorean argument above is direction-blind: in the underlying 4D ontology (\Cref{subsec:ontology-kinematics}) any displacement perpendicular to a link stretches that link, regardless of which ambient direction is perpendicular.
The slice-projected description above treats $u=X^4$ specially because the inside observer has already chosen $X^4$ as the timelike direction.
The same mechanism, applied symmetrically across the four directions, predicts that displacements in the three \emph{spacelike} directions (i.e.\ in-brane kinetic energy / spatial motion) also lengthen the link along the timelike direction --- which is gravitational \emph{time dilation} sourced by stress--energy.
Length contraction (timelike-direction excitation stretches spacelike links) and time dilation (spacelike-direction excitation stretches the timelike link) are therefore two faces of one rule.
The slice-restricted analysis of this subsection captures the first face; a full 4D treatment of the brane action would express both within one expression.

This provides a concrete mechanism for a long-range channel often described informally as ``gravity'' in the present
program: localized transverse gradients can act as sources of slowly varying background strain (contraction/dilation),
and linear wave packets propagating on top of that background experience branch-dependent effective coefficients
(cf.\ \Cref{sec:emergent-relativity}).
In a weak-field, slowly varying regime one may therefore seek a Newtonian-limit reduction in which a scalar potential
built from coarse-grained strain (or elastic energy density) obeys a Poisson-type equation.
Deriving such a reduction quantitatively (and relating the resulting coupling constant to $G$) is an open problem of
the present theory fragment.
